\chapter{Introduction}
\label{ch:introduction}

Decision making under uncertainty is a central challenge in both theoretical
and applied contexts, ranging from financial markets to artificial
intelligence. Understanding how agents adapt their strategies when faced with
noisy or incomplete information provides valuable insight into how learning
and optimization unfold in dynamic environments. In particular, strategies
that rely on probabilistic reasoning---such as Bayesian updating, the Kelly
criterion, and multi-armed bandits---offer different ways of balancing risk
and reward across time.

%---------------------------------------------------------------
\section{The Core Problem: Sticky Priors and Regime Change}
%---------------------------------------------------------------

Standard probability theory, as applied to sequential decision-making,
assumes that the future resembles the past: the statistical properties of
payoffs are stationary, and an agent's accumulated historical data is
informative about what comes next. In equilibrium financial markets, this
assumption is productive. In \emph{regime-changing} environments---markets
experiencing a structural break, a crisis, or a volatility shock---the
assumption breaks catastrophically.

The failure mode is specific. A Bayesian agent that has accumulated 100
observations from a bull market holds a posterior over the payoff probability
that is tightly concentrated near the bull-regime mean. When a bear regime
begins, the agent's new observations from the bear regime must overcome
the inertia of those 100 bull-regime observations before the posterior
shifts materially. We term this the \emph{Sticky Prior Paradox}: the
strength of historical evidence, which is a virtue in stationary environments,
becomes a liability under regime change.

%---------------------------------------------------------------
\section{Research Objectives}
%---------------------------------------------------------------

The primary objective is to develop a simulation framework that acts as a
rigorous testbed for sequential decision-making strategies under both
stationary and non-stationary conditions. Concretely, we evaluate:

\begin{enumerate}
    \item How the exploration versus exploitation dilemma manifests when
    agents must size continuous bet fractions rather than choose among
    discrete actions.

    \item The trade-off between maximising long-term wealth (Full Kelly)
    and maximising the probability of reaching a short-term target while
    avoiding ruin (Fractional Kelly or CPPI).

    \item The robustness---or fragility---of standard Bayesian updating
    frameworks when faced with non-stationary, regime-switching payoffs,
    with Student-$t$ heavy tails.

    \item Whether a Volatility-Augmented HMM can materially reduce
    regime-detection latency, and at what cost in terms of foregone
    growth.

    \item The transaction cost sensitivity of different algorithmic
    strategies, quantifying the ``net-of-fees'' cost of survival.
\end{enumerate}

%---------------------------------------------------------------
\section{Contributions}
%---------------------------------------------------------------

This thesis makes the following contributions:

\begin{enumerate}
    \item \textbf{A unified simulation framework} that evaluates eight
    sequential decision-making agents under Common Random Numbers, ensuring
    statistically valid comparisons across three synthetic and four
    empirical market regimes.

    \item \textbf{Volatility-Augmented HMM (Vol-HMM)}: A two-dimensional
    HMM Bayesian filter that augments return observations with rolling
    volatility and incorporates a CUSUM change-point detector, reducing
    empirical regime-detection lag from 15--20 steps to under 2 steps.

    \item \textbf{Risk-Constrained Kelly (CPPI)}: A CPPI-style drawdown
    floor enforcement that achieves zero ruin in tested scenarios, with
    an explicit quantification of the growth cost of the protection.

    \item \textbf{Transaction Cost Analysis}: Empirical demonstration that
    high-turnover agents incur disproportionate fee drag, motivating
    regime-aware allocation that reduces rebalancing frequency.
\end{enumerate}

%---------------------------------------------------------------
\section{Methodological Overview}
%---------------------------------------------------------------

This work bridges Bayesian statistical modeling with stochastic processes and
computational experimentation. Under the guidance of Professor Shiwei Lan,
the project emphasises Bayesian inference and posterior calibration. Professor
Nicolas Lanchier's expertise on probabilistic foundations and filtration
theory informs the rigorous measure-theoretic setup in
Chapter~\ref{ch:foundations}. By comparing simulation results with
theoretical expectations, the thesis constructs a comprehensive understanding
of adaptive probabilistic decision systems.

%---------------------------------------------------------------
\section{Thesis Outline}
%---------------------------------------------------------------

The remainder of the thesis is organised as follows.
Chapter~\ref{ch:literature} reviews the relevant literature on Kelly
criterion theory, Bayesian bandits, Hidden Markov Models, and portfolio
insurance. Chapter~\ref{ch:foundations} establishes the formal mathematical
framework. Chapter~\ref{ch:methodology} describes the simulation environment,
all agents, and the two proposed methods with their formal guarantees.
Chapter~\ref{ch:experiments} presents experimental results and sensitivity
analyses. Chapter~\ref{ch:conclusion} (the conclusion chapter) synthesises
the findings and outlines future directions.
